\section{An extremal weight and its cross-mass bound}
\label{sec:optimizer}

We first choose vertex weights for which a variational inequality controls
the mass of edges between two sets.  Clique exclusion enters through a
random modification that keeps every weight unchanged in expectation while
making the expected edge mass small.

For a finite simple graph $G$ and a real vector $a=(a_v)_{v\in V(G)}$, put
\[
 M_G(a)=\sum_{\{u,v\}\in E(G)}a_ua_v.
\]
Thus this quadratic expression is defined even when some coordinates are
negative.  For nonnegative weights $w$, write $w(S)=\sum_{v\in S}w_v$ and
\[
 e_w(S,D)=\sum_{u\in S}\sum_{v\in D\cap N_G(u)}w_uw_v.
\]
This is a directed cross mass: an edge with both ends in $S\cap D$ is
counted twice.  All logarithms are natural, and $0\log 0=0$.

\subsection{The extremal weight}

The functional below is the unit-parameter specialization of the
entropy-minus-edge potential used by Davies
\cite[Section~3.3, proof of Theorem~17]{Davies2026}.
We give its elementary maximizing properties here.

Define
\[
 F_G(w)=\sum_{v\in V(G)}w_v(1-\log w_v)-M_G(w),
 \qquad F_G^*=\max_{w\ge0}F_G(w).
\]
The maximum in this definition exists, as the following lemma shows.

\begin{lemma}\label{lem:optimizer}
For every nonempty finite simple graph $G$, the function $F_G$ attains its
maximum.  Every maximizer $w$ has positive coordinates and satisfies
\begin{equation}\label{eq:stationarity}
 \log(1/w_v)=w(N_G(v)) \qquad(v\in V(G)).
\end{equation}
In particular, $w_v\le1$.  For every nonnegative vector $q$,
\begin{equation}\label{eq:optimizer-divergence}
 0\le F_G(w)-F_G(q)
 =\sum_v\left(q_v\log\frac{q_v}{w_v}-q_v+w_v\right)
   +M_G(q-w).
\end{equation}
Here a term with $q_v=0$ in the sum equals $w_v$.
\end{lemma}

\begin{proof}
The continuous function $f(t)=t(1-\log t)$ on $[0,\infty)$ is at most $1$
and tends to $-\infty$ as $t\to\infty$.  Since $M_G(w)\ge0$ for $w\ge0$,
\[
 F_G(w)\le f(\|w\|_\infty)+|V(G)|-1.
\]
Consequently $F_G(w)\to-\infty$ as $\|w\|_\infty\to\infty$, so continuity
gives a maximum.  If a maximizer had $w_v=0$, changing just this coordinate
to $t>0$ would increase $F_G$ by
\[
 t\bigl(1-\log t-w(N_G(v))\bigr)>0
\]
for sufficiently small $t$.  All coordinates are therefore positive, and
the first derivative in each coordinate gives \eqref{eq:stationarity}.
To obtain \eqref{eq:optimizer-divergence}, expand the quadratic expression
$M_G(q)$ around $w$.  Its linear term is
$\sum_v(q_v-w_v)w(N_G(v))$, which cancels the corresponding term using
\eqref{eq:stationarity}.  The inequality follows from maximality.
\end{proof}

For the empty graph on no vertices we set $F_G^*=0$; no positive-coordinate
assertion is needed in that case.

\subsection{Random multipliers with small expected edge mass}

The recursive use of clique-free neighborhoods goes back to the
sparse-subgraph argument of Ajtai, Erd\H{o}s, Koml\'os and
Szemer\'edi~\cite[Section~3]{AEKS1981}. Here the recursion preserves
each vertex weight in expectation. The next lemma applies to arbitrary
positive weights and does not require the maximizing equations.

The construction successively removes neighborhoods from the remaining
induced graph, until every remaining vertex has small neighborhood weight.
Each removed part excludes a smaller clique, and the remainder has small
edge mass. We then retain one part at random and divide its weights by
that part's selection probability.
Every vertex keeps its original expected weight, while edges joining
different parts disappear. Within a selected neighborhood, the same
construction is applied recursively. The multiplier bound below measures
the cost of these successive rescalings.

\begin{lemma}\label{lem:sparse-multipliers}
Let $r\ge2$ be an integer, let $G$ be a finite $K_r$-free graph with positive vertex
weights $w$, and put $s=w(V(G))$.  For every $0<\varepsilon\le1/2$ there
are jointly distributed random variables $(m_v)_{v\in V(G)}$ such that
\begin{equation}\label{eq:sparse-multipliers}
 0\le m_v\le b_r(\varepsilon)
 :=2^{r(r-2)}\varepsilon^{-(r-2)},\qquad
 \mathbb E m_v=1,\qquad
 \mathbb E M_G((w_vm_v)_v)\le\varepsilon s^2.
\end{equation}
Their joint distribution can be taken to have finite support.
\end{lemma}

\begin{proof}
The empty graph on no vertices is immediate.  We induct on $r$, starting
with $r=2$, when there are no edges and $m_v=1$ suffices.  Suppose $r\ge3$
and set $\eta=\varepsilon/4$.

Starting with all vertices, repeatedly choose a vertex whose open
neighborhood in the remaining induced graph has weight greater than
$\eta s$, and remove that neighborhood as a set $V_i$.  These disjoint
sets have weights $s_i=w(V_i)>\eta s$, and each $G[V_i]$ is
$K_{r-1}$-free.  Indeed, a $(r-1)$-clique in such a neighborhood together
with its chosen vertex would form a $K_r$.  When the procedure stops,
the remaining set $R$ satisfies
\begin{equation}\label{eq:remainder-mass}
 M_{G[R]}(w)
 =\frac12\sum_{v\in R}w_vw(N_G(v)\cap R)
 \le\frac12\eta s\,w(R)\le\frac12\eta s^2.
\end{equation}

If no sets $V_i$ were removed, take every $m_v=1$.  Otherwise put
$b=\sum_i s_i>0$ and $p_i=s_i/(2b)$.  Choose $R$ with probability $1/2$
and $V_i$ with probability $p_i$, putting all multipliers outside the
chosen set equal to zero.  On choosing $R$, set $m_v=2$ for $v\in R$.
On choosing $V_i$, use the induction hypothesis there with precision
$\varepsilon/4$, and multiply its output by $1/p_i$.
The option of choosing $R$ is retained even if $R$ is empty; in that
event all multipliers are zero.

Every vertex has mean multiplier one.  Since
$p_i>\eta/2=\varepsilon/8$, each multiplier is at most
\[
 \max\left\{2,\frac8\varepsilon
                    b_{r-1}(\varepsilon/4)\right\}
 = b_r(\varepsilon).
\]
The equality follows directly from the displayed formula for $b_r$.
Only edges within the chosen set contribute.  Therefore
\begin{align*}
 \mathbb E M_G((w_vm_v)_v)
 &\le 2M_{G[R]}(w)
       +\sum_i\frac{\varepsilon s_i^2}{4p_i}\\
 &\le \eta s^2+\frac{\varepsilon}{2}b^2
 \le\frac34\varepsilon s^2.
\end{align*}
The finite construction and the induction hypothesis also give finite
support for the joint law.
\end{proof}

\subsection{The cross-mass estimate}

For $x\ge1$, define
\[
 g_x(s)=
 \begin{cases}
 s\bigl(x+\log^+(1/s)\bigr),&s>0,\\
 0,&s=0,
 \end{cases}
 \qquad \log^+ a=\max\{0,\log a\}.
\]
The function $g_x$ is nondecreasing on $[0,\infty)$.  On $(0,1)$ its
derivative is $x+\log(1/s)-1\ge0$, and on $(1,\infty)$ it is $x$;
the formulas agree continuously at $s=1$ and $s=0$.

\begin{lemma}\label{lem:cross-mass}
Let $r\ge2$ be an integer and let $w$ maximize $F_G$ on a nonempty finite $K_r$-free
graph $G$.  For every $x\ge1$ and all vertex sets $S,D$ with
$w(S),w(D)\le e^x$,
\begin{equation}\label{eq:optimizer-cross-mass}
 e_w(S,D)\le C_r g_x(w(S)),
 \qquad C_r=16r^2.
\end{equation}
The sets $S$ and $D$ need not be disjoint.
\end{lemma}

\begin{proof}
First suppose $S\cap D=\varnothing$, and put $s=w(S)$ and $t=w(D)$.
The case $s=0$ is immediate.  Initially assume $e^{-x}\le s\le e^x$.
We will increase the weights on $S$ and set those on $D$ to zero.
In the variational identity, this creates a negative quadratic term
proportional to $e_w(S,D)$. The random multipliers keep the positive
quadratic term within $S$ small, while their bounded range controls the
vertex terms through a logarithm.
Apply Lemma~\ref{lem:sparse-multipliers} to $G[S]$ with
$\varepsilon=e^{-8x}$, and set $k=e^{4x}$.  Define a nonnegative random
vector by
\[
 q_v=
 \begin{cases}
 (1+km_v)w_v,&v\in S,\\
 0,&v\in D,\\
 w_v,&v\notin S\cup D.
 \end{cases}
\]
The expected quadratic term in \eqref{eq:optimizer-divergence} is at most
\begin{equation}\label{eq:cross-quadratic}
 \mathbb E M_G(q-w)
 \le k^2\varepsilon s^2+\frac12t^2-k e_w(S,D).
\end{equation}
Here the last term uses $\mathbb E m_v=1$; no independence between the
multipliers is required. The choice $k^2\varepsilon=1$ keeps the first
term controlled despite the increase in the weights on $S$.

For $y\ge0$,
$(1+y)\log(1+y)-y\le y\log(1+y)$.  Also
\[
 \log(1+kb_r(e^{-8x}))\le A_rx,
 \qquad
 A_r=4+8(r-2)+(1+r(r-2))\log2.
\]
It follows that the expected sum of vertex terms in
\eqref{eq:optimizer-divergence} is at most $t+kA_rxs$.
Combining this with \eqref{eq:cross-quadratic} and dividing by $k$ gives
\begin{align*}
 e_w(S,D)
 &\le A_rxs+e^{-4x}s^2+e^{-4x}(t+t^2/2)\\
 &\le (A_r+3)xs.
\end{align*}
For the last inequality, use $s,t\le e^x$ and $s\ge e^{-x}$: the three
error terms divided by $s$ are at most $e^{-3x}$, $e^{-2x}$, and
$e^{-x}/2$, respectively, and their sum is at most $3x$.

If instead $0<s<e^{-x}$, apply the estimate just proved with
$x'=\log(1/s)>x$.  Both set weights are at most $e^{x'}$ and
$s=e^{-x'}$, so
\[
 e_w(S,D)\le(A_r+3)x's\le(A_r+3)g_x(s).
\]
Thus the bound with constant $A_r+3$ holds for all disjoint sets.

For arbitrary $S,D$, assign each vertex independently to one of two
classes with equal probabilities.  Let $S'$ be the part of $S$ in the
first class and $D'$ the part of $D$ in the second.  These sets are
disjoint, their weights satisfy the same upper bounds, and
\[
 \mathbb E e_w(S',D')=\frac14e_w(S,D).
\]
There are no loops, so each directed edge in this expectation has two
distinct endpoints.  The disjoint-set bound and monotonicity give
\[
 \frac14e_w(S,D)
 \le (A_r+3)\mathbb E g_x(w(S'))
 \le (A_r+3)g_x(s).
\]
Finally, $A_r+3\le r^2+6r-8\le4r^2$ for $r\ge2$, which proves
\eqref{eq:optimizer-cross-mass}.
\end{proof}

Lemma~\ref{lem:cross-mass} is the only consequence of extremality
needed in the weighted triangle argument.  That argument can therefore
be formulated for any positive weighted graph satisfying
\eqref{eq:optimizer-cross-mass} with some constant $C$.
